%%%%%%%%%%%%%%%%%%%%%%%%%%%%%%%%%%%%%%%%%%%%%%%%%%%%%%%%%%%%%%%%%%%%%
% Based on Tsinghuabeamer(https://github.com/tl3shi/THUBeamer)
% Any question, pls feel free contact me anytime(bhu13@lzu.edu.cn)
% For Lanzhou university beamer
% 编译顺序
% XeLateX->BibTeX->XeLateX->XeLateX
%
%                  ---by suchot
%                  ---2018/9/20
%%%%%%%%%%%%%%%%%%%%%%%%%%%%%%%%%%%%%%%%%%%%%%%%%%%%%%%%%%%%%%%%%%%%%%
% !TEX TS-program = Xelatex
% !TEX encoding = UTF-8 Unicode
\documentclass{beamer}
\usepackage{latexsym,amssymb,amsmath,amsbsy,amsopn,amstext,xcolor,multicol,xeCJK}
\usepackage{graphicx,wrapfig,fancybox,booktabs}
\usepackage{pgf,pgfarrows,pgfnodes,pgfautomata,pgfheaps,pgfshade}
\usepackage{beamerthemeLanzhouUniversity}

\usepackage[hyperref=true,backend=bibtex,sorting=none,backref=true]{biblatex}

\addbibresource{database.bib}
\setbeamertemplate{bibliography item}[text]


\usepackage{algorithm} %%format of the algorithm
\usepackage{algpseudocode}
\algrenewcommand{\algorithmiccomment}[1]{ $//$ #1}
\usepackage{listings}


\setbeamerfont{footnote}{size=\tiny}  %设置页尾字体大小
\renewcommand*{\bibfont}{\footnotesize}



\begin{document}

\title{Lanzhou university Beamer Theme}
\subtitle{subtitle}
\author{Hu Bin胡斌}
\institute{Lanzhou University兰州大学}
%\date{}

\begin{frame}
\titlepage
\begin{figure}[htpb]
  \begin{center}
	\includegraphics[width=0.2\linewidth]{Lanzhou_University_Logo.pdf}
  \end{center}
\end{figure}
\end{frame}

%%多栏目录
%\section*{Table Of Contents}
%\begin{frame}
%    \frametitle{\secname}
%	\begin{multicols}{2}
%    \tableofcontents
%	\end{multicols}
%\end{frame}

%########
%单栏目录
%########
\section*{Table Of Contents}
\begin{frame}
    \frametitle{\secname}
  %  \tableofcontents
	\tableofcontents[hideallsubsections]%隐藏subsection目录
\end{frame}

%########
%在章节前添加Outline
%########
\AtBeginSubsection[] {
\frame<handout:0> {
\frametitle{Outline}
\tableofcontents[current,currentsubsection]
}
}

\section{Section I}
\subsection{Subsection I}
\begin{frame}
\begin{block}{This is a block}
Write something\cite{kocher99} in here \dots
\end{block}
\begin{alertblock}{This is a alertblock}
Write \alert{something} in here \dots
\end{alertblock}
\end{frame}

\subsection{Subsection II }
\begin{frame}
  %\frametitle{\subsecname~}
	Just items:
    \begin{itemize}[<+-| alert@+>]
    \item
   item a
    \item 
   item b
    \item
   item c
    \end{itemize}
\end{frame}

\begin{frame}
  %\frametitle{\subsecname~}
	Just items:
    \begin{itemize}
    \item
   item a
    \item 
   item b
    \item
   item c
    \end{itemize}
\end{frame}


\begin{frame}
    %\frametitle{\subsecname~Goods}
	The enumarate item:
    \begin{enumerate}[<+-| alert@+>]
    \item
	 Item start with number
    \item
     This is item 2
  \end{enumerate}
\end{frame}

\section{section II}
  
\subsection{section I}
\begin{frame}
引用公式：eq.\ref{eq:sbm}
\begin{equation}
\hat{H}=\frac{\epsilon}{2}\hat{\sigma}_{z}-\frac{\Delta}{2}\hat{\sigma}_{x}+\sum_{k}\omega_{k}\hat{b}_{k}^{\dagger}\hat{b}_{k}+\sum_{k}\frac{g_{k}}{2}\hat{\sigma}_{z}(\hat{b}_{k}+\hat{b}_{k}^{\dagger})\label{eq:sbm}
\end{equation}
\end{frame}

\begin{frame}{伪代码}

\scalebox{0.5}{
    \begin{minipage}{1.0\textwidth}
           \begin{algorithm}[H]
              \caption{Wrong algorithm}
              \label{alg:algorithm}
              \begin{algorithmic}[1]
              \Require
              Require 
              \Ensure
             Ensure
              \Function {Tree}{$node_1, node_2$}
                \If{$node_1.bv \cap node_2.bv = \emptyset$}
                  \State \Return{\textbf{False}}
				  \State \Call{Tree}{parameter}
                  \Comment{Comment}
                \Else
                         \ForAll {$child \in node_1.children$}
                         \State \Call{Tree}{$child, node_2$}  \Comment{递归调用}
                         \EndFor
                \EndIf
              \EndFunction
              \end{algorithmic}
              \end{algorithm}
\end{minipage}
}
\end{frame}


\begin{frame}{表格}
引用表格：tab.\ref{tab:Greek}
\begin{table}[H]
\begin{centering}
\begin{tabular}{|c|c|c|c|c|}
\hline
Alpha & Beta & Gamma & Delta & Theta\\
\hline
$\alpha$ & $\beta$ & $\gamma$ & $\delta$ & $\theta$\\
\hline
$A$ & $B$ & $\Gamma$ & $\Delta$ & $\Theta$\\
\hline
\end{tabular}
\par\end{centering}
\protect\caption{希腊字母表\label{tab:Greek}}
\end{table}
\end{frame}

\begin{frame}{三线表}
\begin{table}[htbp]
 \caption{\label{tab:test}三线表}
 \begin{tabular}{lcl}
  \toprule
  。。 & 。。 & 。。 \\
  \midrule
 。。 & 。。 & 。。 \\
 。。 & 。。 & 。。 \\
 。。 & 。。 & 。。 \\
 \bottomrule
 \end{tabular}
\end{table}
\end{frame}


\begin{frame}{页脚}
  Description item:
  \begin{description}%[<+-| alert@+>]
  \item[description] Item starts with description.\footfullcite{kocher99}
  \item[Fermion] This is item 2.\parencite{kocher99}
  \end{description}
\end{frame}
\section{Section III}

\subsection*{bibliography}
    \frame[t,allowframebreaks]{
     % \frametitle{\subsecname}
    \printbibliography
    }

\subsection*{Q\&A}
\begin{frame}
  %\frametitle{\subsecname}
\begin{center}
    {\bf \Huge Q\&A?}\\
\pause
   {\bf \Huge Thank you!}
\end{center}
\end{frame}



\end{document}


%\begin{figure}
%\includegraphics[height=10cm,width=12cm]{a3.eps}
%\caption{}
%\label{a3}
%\end{figure}
